\documentclass[11pt]{article}
\usepackage[margin=1in]{geometry}
\usepackage{amsmath,amssymb,amsthm}
\newtheorem{theorem}{Theorem}
\title{LSPV-2/418: Clause-to-LOCAL Typing and a Paid AND3 Adapter}
\author{PNP SAT-Circuit Reduction Builder}
\date{November 5, 2026, 07:12 PST}
\begin{document}
\maketitle
% LSPV-2 REPORT::418::2026-11-05_0712_PST
\paragraph{Status.} \textbf{PROVED/RESTRICTED} for a specified original-grammar reading that requires a derived factor/equation at a LOCAL input and permits aliased AND3 premises. Finite proof fixtures are \textbf{COMPUTATIONAL}; specialist review pending. Full original Rmix341 grammar, B1/B2/B3, and P versus NP remain \textbf{OPEN}.

\paragraph{IDEA before SEARCH.}
The original \texttt{R\_mix341} rules in LSPV-2/341 (authenticated blob c87acd11) describe LOCAL-CONVERT from one already derived local factor or equation. Proposed CANON2 (LSPV-5/361, blob 7b6e8478) allows LOCAL to consume any clause, including a free source clause, without checking the premise kind. Hypothesis: this overaccepted but sound input typing can be repaired by a separately paid AND3(C,C) which yields a factor exactly equal to C. Falsifiers: aliased premises forbidden, support outside three original variables, false LOCAL implication, corrupt source, invalid factor and intermediate reference.

\paragraph{SEARCH and prior work.}
Read governing 2026-11-02 12:04 five-role handoff first, then issue \#3 through REFEREE363, AGENTS.md, CODEX.md, repo skill, original LSPV-2/341, CANON2/361, CANON2R/362, LSPV-2/405 and /415--417, LSPV-1/418, LSPV-3/423 and LSPV-4/413--414. Read actual GROKPROOF/2026-11-05\_0718\_PROGRESS.md (blob 9c03075) and research/agents/claude/2026-10-28\_1452\_CLAUDE\_TPORT-0002.tex (blob 1f9fc727); neither supplies B3. Repository searches included LOCAL clause parent AND3 alias, LOCAL clause operand adapter, LOCAL clause parent type mismatch. General web search read Haken 1985, Cook--Reckhow 1979, Extended Resolution PHP and Raz--Tzameret Res(linear equations) bibliographic and abstract records. Authenticated GitHub OpenAI searches for LRAT and CNF checker found no relevant native translator; full file read of openai/math/lean/OAI/Algorithms/SubexponentialSAT/ExecCNFNodeSearchCosts.lean (blob 1ced21f4) concerns a specific resource-counted execution. Direct Google-branded search was unavailable; indexed web search substituted. This result neither generates a new hard family nor reclaims established theorems as novel.

\begin{theorem}[Restricted source-authenticated paid conversion]
Fix original Rmix341 sources, unit weights and binary-premise rules. Assume that LOCAL requires a previously derived factor/equation, two AND3 antecedent references may denote the same condition, and all cited clauses with LOCAL support have at most three original variables. Every otherwise legal source-authenticated q-node derivation with u clause-parent LOCAL nodes has a strict-factor-LOCAL derivation of the same final conclusion using
\[
q'\le q+u\le2q.
\]
The independent LSPV-2/417 GF(2)-to-AND3 type adaptation adds at most g distinct paid LOCAL nodes, so both edits compose to give \(q'\le q+u+g\le3q\). For indexed canonical proof records with original input length N bits and fully supplied old evidence length E bits, the rewritten proof has \(E'=O(E+q(N+\log(q+2)))\) bits.
\end{theorem}
\begin{proof}
Each offending clause C has support S of at most three original variables. Introduce one paid AND3 from two references to exactly C. Its truth table on every original assignment is \(C\land C=C\), hence is an equivalent factor \(T_C\) with support S, even for tautological C. Redirect each LOCAL consumer of C to the new factor. The implication check, output and syntactic support union are unchanged; the revised LOCAL now has an already derived factor input. Share each adapter across descendants, and topologically reindex existing node identifiers. A previous LSPV-2/417 affine-to-AND3 adapter creates no new clause-parent LOCAL and therefore does not interfere. Source clauses, paid DEF semantics and final contradiction are unchanged. Reading and rewriting actual evidence has time polynomial in N+E, not necessarily in N+q if arbitrary large payloads are permitted.
\end{proof}

\paragraph{ITERATE, exact proof-DAG test.}
The externally supplied 26-byte DIMACS text is
\begin{verbatim}
p cnf 2 3
1 0
-1 2 0
-2 0
\end{verbatim}
It encodes UNSAT \(F=(x)\land(\neg x\lor y)\land(\neg y)\). The permissive three-node proof applies LOCAL directly to the free clause x, obtaining the x truth-table factor; AND3 of that factor with \(\neg x\lor y\) yields \(x\land y\); AND3 with \(\neg y\) yields the false factor. The literal factor/equation-only interpretation rejects the first inference. The repaired four-node proof adds AND3(x,x) first, then the old LOCAL and the two unchanged consumers. Both derivations are sound; every node is an ancestor of contradiction. Source F has a shorter ordinary RES refutation, so neither exhibited trace is a certified global minimum. The corresponding SAT control \(F'=(x)\land(\neg x\lor y)\land(y)\) has witness \(x=y=1\).

\paragraph{Verification and resource ledger.}
The local Python producer plus independently written restricted checker passed all original-model evaluations, accepted the three-node permissive and four-node strict traces, rejected strict interpretation of the permissive trace, and rejected nine corrupted/malformed or externally changed source controls. Permissive and strict canonical evidence lengths are 598 and 728 bytes, respectively, excluding the external original raw bytes; local truth-table rows are 10 and 12. Independent JavaScript checked four assignments, both labels and five additional corrupted traces (reported 74 assertion groups). The complete immutable proof DAG and a GitHub replay are published separately. This check is for the LOCAL/AND3 fragment, not for every historical rule. The complete peer-owned Python CANON2R checker was inspected through GitHub but was not run verbatim in the local environment. Physical source-byte acquisition, proof reading E, node indexing, output, truth verification and simultaneous memory are additional to paid q; eager implementation space is at most O(N+E) bits, with no FLMV machine transfer.

\paragraph{Outstanding grammar and B3.}
Historical source wording does not say unambiguously whether source clauses are silently promoted to factors or whether two AND3 parent references may alias. Under a permissive reading the extra gate is unnecessary; under the stated strict reading it is sufficient. Neither branch defines private exported LOCAL circuit wires or DEF GF(2) export slots, and neither completes native normalization. Haken's PHP theorem applies to every ordinary shared-DAG Resolution refutation of explicit unary PHP(h+1,h), requiring \(2^{\Omega(h)}\) nodes for formula bit-length polynomial in h. It is not a lower bound on extended resolution, GF(2), full mixed R, or arbitrary SAT machines. Cook ER PHP and Tseitin affine shortcuts are controls. The present proof-to-proof construction does not imply that a polynomial-time SAT decider has polynomial Rmix341 refutations; the weakest bridge
\[
\mathrm{SAT}\in\mathrm P\Longrightarrow
\exists C,k\ \forall N:\overline W_{R_{\mathrm{mix341}}}(N)\le C(N+1)^k
\]
remains \textbf{OPEN}. No proof or counterexample to B3 has been produced.

\paragraph{Five-role handoff.}
LSPV-1: adjudicate the exact original clause/factor LOCAL input type. LSPV-3: determine original AND3 aliased-parent legality. LSPV-4: independently charge E and source-byte acquisition. LSPV-5: freeze explicit selector grammar and check whether direct clause-parent LOCAL is extra. LSPV-2: extend to public LOCAL circuit root and paid DEF wire export. No schedules, enabled states, rules, weights or peer manuscript regions changed.

\begin{thebibliography}{9}
\bibitem{H85} A. Haken, The Intractability of Resolution, Theoretical Computer Science 39 (1985), 297--308.
\bibitem{CR79} S. Cook and R. Reckhow, The Relative Efficiency of Propositional Proof Systems, J. Symbolic Logic 44 (1979), 36--50.
\bibitem{Buss87} S. Buss, Polynomial Size Proofs of the Propositional Pigeonhole Principle, J. Symbolic Logic 52 (1987).
\bibitem{RT08} R. Raz and I. Tzameret, Resolution over Linear Equations and Multilinear Proofs, Annals Pure Appl. Logic 155 (2008), 194--224.
\end{thebibliography}
\end{document}
% END_OF_VERIFIED_REPORT::LSPV-2::418
